% ============================================================
%  Poul Martin Møller, En Dansk Students Eventyr (1824)
%  First published in Dansk Læsebog for Almue-Skolerne
%  (1829 expanded version). Standard text: Efterladte Skrifter,
%  vol. 3, ed. F.C. Olsen (Copenhagen: C.A. Reitzel, 1843).
%
%  TRANSCRIPTION (Danish)
%  Completed: The Licentiate passage
% ============================================================

\documentclass[12pt]{article}
\usepackage[utf8]{inputenc}
\usepackage[T1]{fontenc}
\usepackage[danish]{babel}
\usepackage[protrusion=true,expansion=true]{microtype}
\usepackage[margin=1.4in]{geometry}
\usepackage{setspace}
\usepackage{libertinus}
\usepackage{libertinust1math}
\usepackage{fancyhdr}
\usepackage[colorlinks=true,linkcolor=black,urlcolor=black,
  pdftitle={En Dansk Students Eventyr (uddrag)},
  pdfauthor={Poul Martin Møller}]{hyperref}
\setlength{\headheight}{15pt}
\pagestyle{fancy}
\fancyhf{}
\fancyhead[LE,RO]{\thepage}
\fancyhead[RE]{\textit{En Dansk Students Eventyr}}
\fancyhead[LO]{\textit{Licentiaten}}
\setlength{\parindent}{1.5em}
\setlength{\parskip}{0pt}
\onehalfspacing

\title{\textbf{En Dansk Students Eventyr}\\[0.5em]
  \large Uddrag: Licentiaten}
\author{Poul Martin Møller}
\date{Første gang trykt 1824; revideret udgave i
  \textit{Efterladte Skrifter}, bd. 3,\\
  red. F.C. Olsen. København: C.A. Reitzel, 1843}

\begin{document}
\maketitle

\section*{Licentiaten}

\noindent
For tre Aar siden havde Familien paa Ravnshøi og adskillige af deres
velhavende Naboer, der ei have nogen offentlig beskikket Medicus i
deres Nabolaug, forenet sig om at holde sig en Læge for egen Regning.
Jeg besluttede da at gribe denne Leilighed, for at bringe min kjære
Frænde til et Slags Virksomhed. Thi selv er han alt for lam paa
Villien til at udføre nogen Beslutning. Jeg gav ham et Brev i Lommen,
satte ham paa en Fragtvogn og troede, Sagen var fuldkommen i sin
Rigtighed. Men da jeg et Aars Tid efter begav mig derud for at see,
hvorledes han skikkede sig i sin nye Stilling, og i denne Hensigt lod
mig føre til Stedets Læge, fandt jeg til min store Forundring, at en
anden Medicus øvede sin Kunst der i Egnen. Nu vilde jeg igjen reise
til Staden, for om muligt der at opspørge Licentiaten. Paa Veien
overnattede jeg i Møllekroen og saae der, mens jeg i Aftenskumringen
sad med den lystige Vært og betragtede Kakkelovnsilden, en lang,
kroget Skikkelse liste sig langs med Væggen og forsvinde gjennem en af
Dørene. Denne Skikkelse vil jeg beskrive; thi med det Samme beskriver
jeg min Fætter. Han er halvfjerde Alen lang, og af Skamfuldhed over
denne overnaturlige Længde, som ved hans Smekkerhed endnu mere bliver
iøinefaldende, bøier han sin Ryg saameget, at hans Hoved falder under
et almindeligt Menneskes Synskreds. Naturen har lettet ham denne
Stilling ved at forsyne hans Rygrad med to Led, saaledes at den, i
Særdeleshed naar han sidder, udgjør tre Linier, der danne to stumpe
Vinkler med hinanden. Forresten er der Intet at mærke ved hans Ydre,
naar man undtager en Klumpfod, som gjør ham det nødvendigt, til Brug
for sin høire Fod at holde sig en Snørestøvle.

--- Hvad var det for et Menneske? spurgte jeg Værten.

--- Det er een af mine Gjæster, gav han til Svar; han kom for et Aar
siden til mit Huus og lod til, at have stort Hastværk: thi han
bestilte Thevand til næste Morgen Kl.~4, for ret betids at fortsætte
sin Reise. Men den næste Dag lod han Vognen afsige. Dette gjentog han
siden hver Dag i det første halve Aars Tid; men nu lader det til, at
han har slaaet sig til Ro her i mit Huus. Det er et godt Stykke Karl,
som aldrig kommer nogen Sjæl for nær; men jeg troer, sandt at sige,
han er ikke vel forvaret. Hans Lærdom har gjort ham forrykt i Hovedet.

--- Valg Deres Ord forsigtig, sagde jeg, thi han er mit kjødelige
Søskendebarn paa mødrene Side.

--- Ja, sagde han; saa maa De jo selv bedst vide, hvordan han har
det. Han er forfærdelig folkesky. Igaar lod han en Stige sætte til
sit Kammervindue, for at kunne komme ned i Haven og gjøre sig sin
sædvanlige Motion, uden at gaae igjennem Gjæstestuen. Sommetider
ligger han en heel Dag i Dvale, sommetider spanker han op og ned ad
Gulvet og giver sig til at præke for sig selv. Men det er jo gandske
overflødigt, at jeg fortæller Sligt til Dem. De maa jo selv bedst
kjende hans Underligheder. Han gjør os, som sagt, ingen Bryderier. Vi
ere blevne saa vante til at see ham humpe omkring her i Huset, at jeg
endogsaa troer, vi vilde savne ham, dersom han reiste sin Vei. Min
Datter Marie passer meget omhyggelig paa, at han i betimelig Tid faaer
sin Mad og Drikke, ellers troer jeg, han af Glemsomhed sultede reent
ihjel. Jeg siger tidt nok til hende: Lad os engang see, om ikke hans
Natur selv kan minde ham. Han er dog saa lang, at han maa kunne kræve
sig selv. Men hun er altfor blødhjertet til at anstille Forsøget. ---

Jeg sad oppe til langt ud paa Natten for at skrive Breve, og hørte
ham imidlertid bestandig pusle paa sit Værelse. Snart phantaserede han
nogle faa Minuter paa sin Fløite, snart spadserede han taus og med
stærke Skridt omkring paa sit Gulv. Jeg gik ud i Haven, for at kige
gjennem hans Vinduer. Lyset paa hans Bord var brændt langt ned i
Pladen, da han i sine overjordiske Tanker havde forglemt at skyde det
i Veiret. Luen steeg og sank i Blikrøret, der var gandske blaat af
Heden, og den smeltede Talg udbredte sin Stank i hele Værelset. Halv
liggende, halv siddende befandt han sig i en meget lav Lænestol, med
sine lange Been udstrakte paa Gulvet. Han var fordybet i Betragtninger
og stirrede stivt hen for sig igjennem sine Briller. Begge hans
Hænder, der endnu holdt fast paa Fløiten, hvilte aldeles ubevægede paa
hans Skjød. Jeg betragtede ham flere Minuter, uden at han gjorde Mine
til at røre noget Lem. Naar jeg ikke havde seet det brændende Blik
igjennem Brillerne, vilde jeg have taget det for en heldig Voxafstøbning
af min Fætter Licentiaten. Tilsidst raabte jeg til ham gjennem det
aabne Vindue: God Aften, Fætter! --- Han foer sammen uden at bringes
ud af sin forstenede Positur. Nu krøb jeg ved Hjelp af den nylig
anbragte Stige gjennem Vinduet ind til ham, og bragte ham ved mine
Nævers Kraft ud af sin Dvale. Han reiste sig op med endeel forlegne
Bukninger og fremstammede nogle ufuldbaarne Complimenter, hvormed han
yttrede sin Taksigelse, fordi jeg saa gjerne paatog mig den Uleilighed,
at hjelpe lidt paa hans Leveplan, og gjorde mig endeel Undskyldninger,
fordi han ikke havde fulgt den Bane, jeg havde foreskrevet ham. Hans
Tanker skride saa langsomt, at han undertiden mangler Aandsnærværelse
til at svare paa det eenfoldigste Spørgsmaal. I Særdeleshed er han
uforstaaelig, naar han nylig har været hentaget af sine metaphysiske
Phantasier. Da hersker der en saadan Skilsmisse imellem hans Ord og
Mening, at han ei er i Stand til at fremføre nogen fuldstændig
Sætning; men alle hans Tanker komme dødfødte til Verden. Da nogen
Ligevægt var tilveiebragt i hans adspredte Sind, aflagde han sit
Skriftemaal, der, med Forbigaaelse af hans hyppige Pauser og mislykkede
Perioder, lød saaledes:

--- Dagen efterat jeg var kommet hid, var det mit faste Forsæt at
begive mig til Ravnshøi; men det var mig umuligt at finde min
Lommesax, og du vilde dog neppe forlange, jeg skulde fremstille mig
for et Herskab med uklippede Negle. Da jeg den følgende Dag greb en
Brand paa Skorstenen, for at tænde min Pibe, brændte jeg mig en stor
Vable paa den høire Pegefinger. Med en saadan legemlig Smerte begriber
du vel, jeg manglede det frie Hovedomløb, som udkræves, naar man skal
indføre sig selv med en passende Hilsen hos en virkelig Justitsraadinde.
Da Møllerkonen den tredie Dag, efter min Anviisning, skulde forbinde
min syge Haand, og med en tændt Praas vilde brænde Sytraaden over,
dryppede syv store Talgpletter paa min sorte Kjole. Da den første
Draabe var faldet, forudsaae jeg, at der vilde følge flere paa; men
jeg var saa betaget af Taknemmelighed over hendes Omhu, at jeg ikke
havde Mod til at gjøre hende opmærksom derpaa. Nu gik lang Tid hen,
inden jeg fik anskaffet noget Terpentin til at tage disse Pletter af
med. Imidlertid vare mine skjønneste Kraver smudsede, og endeel Uger
forløb, inden jeg fik dem besørgede til en Vaskerkone. Med saadanne
smaa Uheld, som en ond Dæmon sendte mig paa Halsen, forløb et Aars
Tid, uden at det var mig muligt at komme til Ravnshøi, som blot ligger
en halv Miil fra denne Kro.

--- Men min Gud, udbrød jeg, hvorledes er det tænkeligt, at du ei paa
en eneste af de 365 Dage fik det Heltemod, til Trods for alle
Hindringer at gjøre det herkuliske Tag!

--- Saaledes, svarede han, taler bestandig du og Enhver, der har det
Held, at kunne handle uden at tænke over sit Liv. Naar jeg havde
kunnet indsee nogen tilstrækkelig Grund til, at Reisen snarere skulde
skee paa een Dag end paa en anden, da vilde den have gaaet for sig.
Men Ulykken bestaaer egentlig deri, at jeg altfor tydelig indseer den
Sandhed, at enhver Gjerning uden stor Skade kan opsættes en Dag.

--- Men deri tager du just feil, Fætter, gav jeg til Svar; ingen Ting
i Verden, som kan udføres idag, bør udsættes til imorgen.

--- Negter du mig det, tog han atter Ordet, da maa du nødvendig
indrømme mig, at den muligviis ligesaa godt kan udføres et Minut
eller et Secund senere; og denne forbistrede Mulighed varer ved,
saalænge man drager Aande. Ingen kan bevise mig at en Opsættelse, der
varer en Secund eller Terts, skulde have fordærvelige Følger, og fordi
Ingen kan bevise mig det, derfor gaaer mit Liv reent i Staa. Min
uendelige Grandsken derover gjør, at jeg Intet udretter. Fremdeles
kommer jeg til at tænke paa mine Tanker derover, ja jeg tænker over,
at jeg tænker derover, og deler mig selv i en uendelig tilbageskridende
Rad af Jeg'er, der betragte hinanden. Jeg veed ikke, hvilket Jeg der
skal standses ved som det egentlige, og i det Øieblik jeg standser ved
eet, er det jo igjen et Jeg, der standser derved. Jeg bliver ør og
betaget af Svimmelhed, som om jeg stirrede ned i en bundløs Afgrund,
og Tænkningen endes med, at jeg føler en rædsom Hovedpine.

--- Kjære Fætter, sagde jeg rolig, med slige Grublerier vil du pine
dig ligesaa forgjeves, som om du endevendte din Goliathskrop, og
forsøgte at spadsere paa Hovedet med dine lange Been i Veiret. Jeg
kan slet ikke staae dig bi med at sortere dine mange Jeg'er. Det
ligger aldeles uden for min Virkekreds, og jeg maatte enten være eller
blive ligesaa gal som du, hvis jeg indlod mig i dine overmenneskelige
Drømmerier. Min Viis er at holde mig til de haandgribelige Ting og
vandre paa Forstandens brede Landevei; derfor komme mine Jeg'er aldrig
i Ulave. Forresten seer jeg nok, at Forseelsen er min. Jeg burde have
ført dig lige til Maalet, aabnet Justitsraadindens Døre, og stødt dig
ind i Stuen saaledes, at du havde seet dig nødsaget til at gribe det
første det bedste Jeg, for at byde Fruen og hendes gode Familie god
Dag. Men siig mig nu, hvad du har taget dig for her i et udslaget Aar?

--- Jeg har, sagde han, samlet Materialier til mit store Værk over
Menneskets physiske og intellectuelle Natur.

--- Aa, svarede jeg, tænker du paa det endnu? For 10 Aar siden havde
du den samme fixe Idee. Det maa blive ret et grundigt Stykke Arbeid.
Jeg troede, du studerede Mineralogie, siden jeg saae den store Mængde
Stene her paa dit Gulv.

--- De sindes her alene for Værkets Skyld, sagde han. Jeg har bemærket,
at Tankerne glide bedst fra Pennen, naar man besidder aldeles
hensigtsmæssige Skriveredskaber. For at faae gode Penne, maa man have
gode Knive, og for at skjærpe dem, maa man besidde ret udsøgte
Slibestene. Derfor har jeg, inden jeg gik videre, forskrevet adskillige
Steenarter, og rigtig nok, som du seer, maattet trænge temmelig dybt
ind i Mineralogien som et forberedende Studium.

--- Herpaa gav jeg til Svar: Dette synes at være ret en grundig
Fremgangsmaade, naar du hertil føier et flittigt Studium af
Ornithologien, for at løse den Opgave, hvilke Fugles Fjær der yde de
bedste Penneposer, og tillige opoffrer nogle Aar til Pædagogiken, for
at udfinde, hvorledes en Dreng bedst opdrages, for engang i Tiden at
blive en dygtig og retsindig Knivsmed.

--- Egentlig talt var det slet ikke for vidt drevet, sagde han. Du
spotter, fordi du ikke forstaaer mig. Disse Forsøg paa at lette den
aandelige Virksomhed ved Materiens Beqvemhed afgive jo tillige Bidrag
til Værket selv, ligesom ogsaa Betragtningen over, hvorledes jeg
bragtes til Undersøgelsen, udgjør et nyt Capitel. Fremdeles ---

--- Ja, stop nu, Fætter, afbrød jeg ham, ellers løbe vore Tanker atter
Surr, ligesom da vi kom i Vilderede med vore Jeg'er. Man bør ikke roe
saa langt ud, at man taber Landet af Sigte.

For at rede mig selv og min Fætter ud af det filtrede Hjernespind,
hvori han gjerne saaledes indvikler os, at jeg med min jydske Forstand
tilsidst veed hverken ud eller ind, gav jeg mig til at undersøge, hvad
Forfatning hans Klæder befandt sig i; thi han er i enhver praktisk
Dont aldeles raadvild og hjelpeløs. Jeg erfarede da, at hans Garderobe
ved den menneskekjærlige Maries Bistand var bragt i bedre Stand, end
jeg havde troet. Medens jeg syslede hermed, blev jeg vaer, at han,
ved at holde Frakken opknappet, gav sit Liv et bredere Gjennemsnit,
for at spærre mig Udsigten til en vis Krog. Naar jeg vilde kaste mit
Øie derhen, trippede han omkring foran min Synslinie som en Gaas, der
med udspilede Vinger søger at forsvare sit nys udrugede Kuld mod
Hønsepigens Hænder; men hun sparker den kortsynede Gaas tilside med
sine Træskoe og tæller alle Gæslinger i sit Sold, for at bringe dem i
Nattely. Saaledes greb jeg, der havde mærket Uraad, den fortvivlede
Licentiatus ved Armen og sagde: Gaae tilside, Fætter, jeg maa see,
hvad det er for Contrabande, du har skjult her i Krogen. Her opdagede
jeg en Stabel skimlede Kraver og Halsklude, som han meget omhyggeligen
havde vidst at skjule for Maries Speiderblik. Med en Lineal rørte jeg
op i det ilde tilredte Linned, og forstyrtede derved en heel Slægt
Ørentvister i deres huuslige Lyksalighed. Den flygtende Sværm af
Insekter undslap ved Hjelp af sine mangfoldige Been og smidige Kroppe
for største Delen mine morderiske Forfølgelser; men jeg udbrød i Harme:

--- Hjertenskjære Fætter! Hvorfor i al Verden skjuler du din
Skrøbelighed for mig? At du dølger din Uhumskhed for den smukke
Møllerdatter, kan jeg finde forklarligt; men at du vil undsee dig for
din kjødelige Fætter, der saa tidt har seet dig i din hele Nøgenhed,
det gaaer dog altfor vidt.

--- Bertel, svarede han, jeg forstaaer det neppe selv. Det er et
mærkeligt psychologisk Problem. Det gaaer mig som en dødelig Syg, der
ei engang har Mod til at aabenbare Lægen alle sine farlige Tilfælde.
Saaledes deler Mennesket ved mange Leiligheder sig selv i to Personer,
af hvilke den ene søger at føre den anden bag Lyset, imens en tredie,
der i Grunden er den samme som de to andre, forundrer sig høilig over
denne Confusion. Kort sagt, Tænkningen bliver dramatisk og spiller i
Stilhed de meest forviklede Intriger med sig selv og for sig selv; men
Tilskueren bliver bestandig paany Skuespiller. Derom anbringer jeg et
langt Capitel i mit Værk.

\end{document}
